\section*{الفصل \ref{ch:b2:ligand-field} --- نظرية مجال الربائط واللون}

\begin{solution}{exo:b2:ligand-field:1}
$d^3$: $t_{2g}^3$، وCFSE $-1.2\Delta_o$. $d^8$: $t_{2g}^6e_g^2$، وCFSE $-2.4 + 1.2 = -1.2\Delta_o$.
\end{solution}

\begin{solution}{exo:b2:ligand-field:2}
\ce{[Fe(H2O)6]^2+}، $d^6$ \omterm{def:b2:ligand-field:spin}{عالي السبين} $t_{2g}^4e_g^2$: أربعة. \ce{[Fe(CN)6]^4-}، \omterm{def:b2:ligand-field:spin}{منخفض السبين}
$t_{2g}^6$: لا شيء (\omterm{def:b2:diatomic-mos:magnetism}{دايامغناطيسي}).
\end{solution}

\begin{solution}{exo:b2:ligand-field:3}
الطول 600 nm ضوء برتقالي؛ فيبدو المعقد أزرق (أزرق مخضرًا).
\end{solution}

\begin{solution}{exo:b2:ligand-field:4}
$\ce{Cl-} < \ce{H2O} < \ce{NH3} < \ce{CN-}$.
\end{solution}

\begin{solution}{exo:b2:ligand-field:5}
الماء: $\Delta_o < P$، سبين مرتفع، $t_{2g}^3e_g^2$، خمسة إلكترونات منفردة. السيانيد: $\Delta_o > P$،
سبين منخفض، $t_{2g}^5$، إلكترون منفرد واحد.
\end{solution}

\begin{solution}{exo:b2:ligand-field:6}
في المعقد الرباعي الأوجه \ce{[CoCl4]^2-} يبلغ الانشطار نحو $4/9$ من انشطار معقد ثماني الأوجه، 
والكلوريد ربيطة أضعف من الماء: فينتقل نطاق d--d إلى أطوال موجية أطول (من البرتقالي
إلى الأحمر)، ويبدو المعقد أزرق. والمعقدات الرباعية الأوجه تمتص أيضًا بشدة أكبر،
ومن هنا اللون الداكن.
\end{solution}

\begin{solution}{exo:b2:ligand-field:7}
\ce{Ni^2+} توزيعه $d^8$. في المربع المستوي: تحمل المدارات الأربعة السفلى الإلكترونات الثمانية أزواجًا،
ويبقى $\mathrm d_{x^2-y^2}$ فارغًا: \omterm{def:b2:diatomic-mos:magnetism}{دايامغناطيسي}. وفي رباعي الأوجه: $e^4t_2^4$، إلكترونان منفردان
في المجموعة $t_2$.
\end{solution}

\begin{solution}{exo:b2:ligand-field:8}
$\tilde\nu = 1/(500 \times 10^{-7}\,\text{cm}) = \qty{20000}{cm^{-1}}$؛ $\Delta_o =
0.1196/(500 \times 10^{-9}) = \qty{2.39e5}{J/mol}$، أي \qty{239}{kJ/mol}.
\end{solution}

\begin{solution}{exo:b2:ligand-field:9}
\ce{Zn^2+} توزيعه $d^{10}$ و\ce{Sc^3+} توزيعه $d^0$: فلا \omterm{def:b2:ligand-field:d-d}{انتقال d--d} ممكن (لا مستوى d فارغ
في \ce{Zn^2+}، ولا إلكترون d في \ce{Sc^3+})، فلا امتصاص في المرئي.
\end{solution}

\begin{solution}{exo:b2:ligand-field:10}
Co(III) $d^6$ + ستة أزواج حرة للأمونيا: 18 إلكترونًا. اثنا عشر منها تملأ \omterm{def:b2:diatomic-mos:bonding}{المدارات الرابطة} الستة،
وستة تملأ المدارات $t_{2g}$ غير الرابطة (فالأمونيا تعطي مع الكوبالت(III) مجالًا قويًا بما يكفي
لسبين منخفض)؛ و$e_g^*$ فارغ. وكل الإلكترونات مقترنة: \omterm{def:b2:diatomic-mos:magnetism}{دايامغناطيسي}.
\end{solution}

\begin{solution}{exo:b2:ligand-field:11}
\ce{CO} مانح $\sigma$ جيد، وهو قبل كل شيء مستقبل $\pi$: فمداراته $\pi^*$ الفارغة
تخفض المستوى $t_{2g}$ \omterm{def:b2:ligand-field:pi-ligands}{بالتبرع العكسي}، فتوسّع $\Delta_o$. أما \ce{F-} فمانح $\pi$، يرفع
$t_{2g}$ ويضيّق $\Delta_o$. ونموذج الشحنات النقطية لا يرى إلا الشحنة فيخطئ في
الترتيب.
\end{solution}

\begin{solution}{exo:b2:ligand-field:12}
لولا المجال، لاتبعت أنتالبيات الإماهة خطًا منتظمًا (إذ تصغر الأيونات على طول
السلسلة). والماء يعطي معقدات مائية \omterm{def:b2:ligand-field:spin}{عالية السبين} طاقة استقرار مجالها البلوري معدومة من أجل $d^0$، $d^5$، $d^{10}$
وأكبر ما تكون من أجل $d^3$ و$d^8$: فيضيف الاستقرار الإضافي حدبتين، قمتاهما قرب
\ce{V^2+} و\ce{Ni^2+}، فوق الخط المار بالأيونات \ce{Ca^2+}، \ce{Mn^2+}، \ce{Zn^2+}.
\end{solution}

\begin{solution}{pb:b2:ligand-field:1}
\textbf{1.} $d^3$.
\textbf{2.} $t_{2g}^3$، إلكترون واحد في كل من المدارات الثلاثة، بسبينات متوازية.
\textbf{3.} لا: فالإلكترونات الثلاثة تتسع في $t_{2g}$ دون اقتران.
\textbf{4.} $-1.2\Delta_o$.
\textbf{5.} ثلاثة.
\textbf{6.} \omterm{def:b2:ligand-field:cfse}{طاقة استقرار المجال البلوري} كبيرة ومدارات $e_g^*$، المتجهة نحو الربائط،
فارغة: فنزع ربيطة أو إضافتها يكلف جزءًا كبيرًا من ذلك الاستقرار.
\textbf{7.} ستة أيونات أكسيد عند رؤوس ثماني أوجه مشوَّه قليلًا.
\textbf{8.} $1/(556 \times 10^{-7}) = \qty{17990}{cm^{-1}}$، أي نحو \qty{1.80e4}{cm^{-1}}.
\textbf{9.} $0.1196/(556 \times 10^{-9}) = \qty{2.15e5}{J/mol}$: أي \qty{215}{kJ/mol}.
\textbf{10.} الأصفر المخضر (556 nm) والبنفسجي (405 nm).
\textbf{11.} الأحمر، مع قليل من الأزرق: الأحمر الداكن للياقوت.
\textbf{12.} لأيون $d^3$ أكثر من توزيع مثار واحد فيه إلكترون في $e_g$؛
والتنافر بين الإلكترونات يعطيها طاقات مختلفة، ومن هنا النطاقات المتعددة (تُعالج في
مجلد السنة الثالثة).
\textbf{13.} $0.1196/(620 \times 10^{-9}) = \qty{1.93e5}{J/mol}$، أي \qty{193}{kJ/mol}
(\qty{16130}{cm^{-1}}).
\textbf{14.} الياقوت: فانشطاره أكبر.
\textbf{15.} نحو البرتقالي المحمر: فالأحمر يُمتص الآن، ويَنفذ الأخضر (مع بعض الأزرق).
\textbf{16.} في البيريل تكون مواقع الكروم أكبر قليلًا وأيونات الأكسيد أبعد،
والمحيط مختلف: مجال أضعف.
\textbf{17.} لا: فأيون $d^3$ يحتفظ بثلاثة إلكترونات منفردة في أي مجال ثماني الأوجه.
\textbf{18.} $d^9$.
\textbf{19.} يمتص المعقد المائي في الأحمر والبرتقالي \omterm{def:b2:ligand-field:d-d}{بانتقال d--d}: فيكون
الضوء النافذ أزرق.
\textbf{20.} يصير أبيض: فبدون ربائط الماء حول النحاس (تحل الكبريتات
محلها)، يكون المجال أضعف وينتقل النطاق إلى تحت الأحمر.
\textbf{21.} يرتبط الماء بالنحاس من جديد ويتكوّن المعقد المائي الأزرق.
\textbf{22.} عند طول موجة أقصر: فالأمونيا، الواقعة فوق الماء في السلسلة، تعطي
انشطارًا أكبر (الأزرق البنفسجي الداكن لمعقد الأمين).
\textbf{23.} لا: فالتوزيع $t_{2g}^6e_g^3$ هو التوزيع الوحيد.
\textbf{24.} نعم: إلكترون منفرد واحد.
\textbf{25.} $\boldsymbol{\Delta_o \approx \qty{1.8e4}{cm^{-1}}}$، أي نحو
$\boldsymbol{\qty{215}{kJ/mol}}$.
\end{solution}
